\documentclass[11pt]{exam}
\usepackage[empty]{fullpage}
\usepackage{hyperref}
\usepackage{amsmath,amssymb}
\usepackage{color, soul}

\begin{document}
\subsection*{Math 222 - Project 4 \hfill Due Friday, April 3}

\begin{questions}

\question The file 
{\footnotesize
\url{https://bclins.github.io/spring26/math222/CalculusGrades.csv}
}
contains grade data from my calculus classes (Math 141, 142, and 242) from 2013 to 2020.  

\begin{parts}
\item Make side-by-side boxplots to show the grade distribution for each of the three courses. 
\item Which course looks like it has the highest grades?  Which has the lowest grades? 
\item Use one-way ANOVA to see if there is significant evidence of a difference in average grades for the three calculus courses.  Clearly explain your conclusions.
\part Careful analyze whether the conditions for an ANOVA F-test are satisfied.  Include any additional graphs and tables that you need, and clearly explain your findings. 
\part What are the results if you just use a two-sample t-test to compare MATH 141 grades versus MATH 242?  Is there a significant difference?  Is that a big deal?  Explain. 
\end{parts}

\question The link {\footnotesize
\url{http://people.hsc.edu/faculty-staff/blins/classes/spring17/math222/data/bat10.txt}}
contains data from 327 Major League baseball players from one season.  The variable \verb|HR| is the number of home runs hit by each player. 
\begin{parts}
\part Make a graph to visualize how the number of home runs differs for the four different types of position (catchers, designated hitters, in-field, and out-field). 
\part Does it appear that there are significant differences in the average number of home runs for the different positions? Do an ANOVA test and explain your results. 
\part Use the Tukey's Honest Significant Differences command in R to make confidence intervals for the differences between the average number of home runs for each pair of positions.  Which pairs of positions are clearly better?  How much better are they?  
\part Careful analyze whether the conditions for an ANOVA F-test are satisfied.  Include any additional graphs and tables that you need, and clearly explain your findings. 
\end{parts}


%\question Dansinger, Griffith, Gleason et al. (2005) report on a randomized, comparative experiment in which 160
%subjects were randomly assigned to one of four popular diet plans: Atkins, Ornish, Weight Watchers,
%and Zone (40 subjects per diet). These subjects were recruited through newspaper and television
%advertisements in the greater Boston area; all were overweight or obese with body mass index values
%between 27 and 42. Among the variables measured were
%\begin{itemize}
%\item Which diet the subject was assigned to
%\item Whether or not the subject completed the twelve-month study
%\item The subject’s weight loss after two months, six months, and twelve months (in kilograms, with a negative value indicating weight gain)
%\item The degree to which the subject adhered to the assigned diet, taken as the average of 12 monthly
%ratings, each on a 1-10 scale (with 1 indicating complete non-adherence and 10 indicating full
%adherence)
%\end{itemize}
%Data for the 93 subjects who completed the 12-month study are in the file \href{http://people.hsc.edu/faculty-staff/blins/classes/spring16/math222/data/ComparingDiets.csv}{ComparingDiets.csv}.
%Some of the questions that the researchers studied are:
%\begin{parts}
%\item Do the average weight losses after 12 months differ significantly across the four diet plans?
%\item Is there a significant difference in the completion/dropout rates across the four diet plans?
%\item \hl{Is there a significant positive association between a subject’s adherence level and his/her amount of weight loss?}
%\item Is there strong evidence that dieters actually tend to lose weight on one of these popular diet plans?
%\end{parts}
%For each of these research questions, first identify the explanatory variable and the response variable,
%and classify each as categorical or quantitative. Then use graphical and numerical summaries to
%investigate the question, and summarize your findings. Next, identify the inference technique that can
%be used to address the question, and apply that technique. Be sure to include all aspects of the
%procedure, including a check of its technical conditions. Finally, summarize your conclusions for each
%question. Write a paragraph summarizing your findings from these four analyses. [Hint: To determine
%the completion rate for each diet, count how many of the 93 subjects who completed the study are in
%each diet group and compare those counts to the 40 that were originally assigned to each diet.]
\end{questions}
\end{document}
